% GENERATED FILE. Edit canonical lessons, then run npm run generate:books.
% canonical-source-hash: fnv1a64:0994c080336493bb
% canonical-lessons: ZH-C336-gaobie, ZH-C336-huansong, ZH-W336-song-send, ZH-C336-songke, ZH-C336-songli, ZH-C336-woshou

\chapter{Seeing Friends Off}
\label{ch:zh-seeing-friends-off}
\hlchaptermodality{336}

\begin{chapteropening}
\textbf{By the end of this chapter:} \emph{I can say goodbye, give someone a send-off, see a guest to the door, give a present and shake hands.}

Complete the last lesson of chapter 336: I can say goodbye, give someone a send-off, see a guest to the door, give a present and shake hands.
\end{chapteropening}

\section[gàobié]{gàobié — to say goodbye; to take leave}

\label{lesson:ZH-C336-gaobie}



\begin{quote}
Before the new one: say the Mandarin for a riddle, then the Mandarin for a string of firecrackers.
\end{quote}

\subsection*{You'll want to know: gàobié}
\emph{gàobié} — \textquotedblleft{}to say goodbye; to take leave\textquotedblright{}.

\begin{quote}
\emph{tā hé dàjiā gàobié yǐhòu jiù shàng chē le} — After saying goodbye to everyone he got on board.
\end{quote}

Its characters come later; for now it is a word for the ear.

\begin{cousinweb}
\emph{gàobié} is built from \emph{gào} 'to tell' and \emph{bié} 'to part', both inherited from Middle Chinese, \emph{kawH} and \emph{bjet} (Guangyun, in Baxter's transcription), which glosses the first 'to report' and the second 'to differ' and 'to part'.
\end{cousinweb}

\begin{grammarlens}[title={What you've built}]
One more word for telling what happened.
\end{grammarlens}

\subsection*{Your turn}
\begin{itemize}
\raggedright
  \item \emph{Say it:} \emph{gàobié}
  \item \emph{Say it:} \emph{gàobié}, once more
  \item \emph{Say it:} \emph{tā hé dàjiā gàobié yǐhòu jiù shàng chē le} — After saying goodbye to everyone he got on board.
  \item \emph{Recall:} say the Mandarin for a riddle, then the Mandarin for a string of firecrackers, then say \emph{gàobié} again
\end{itemize}

\begin{tcolorbox}[breakable,colback=teal!4,colframe=teal!35!black,title={Before you move on}]
What does gàobié mean? (\textquotedblleft{}To say goodbye; to take leave\textquotedblright{}.) Say it once more.
\end{tcolorbox}
\section[\texorpdfstring{\zh{欢送} (huānsòng)}{欢送 (huānsòng)}]{\zh{欢送} (huānsòng) — to give a send-off}

\label{lesson:ZH-C336-huansong}



\begin{quote}
Before the new one: say the Mandarin for a string of firecrackers, then the Mandarin for to take leave of someone.
\end{quote}

\subsection*{You'll want to know: \zh{欢送}}
\textbf{\zh{欢送}} — \emph{huānsòng} — \textquotedblleft{}to give a send-off\textquotedblright{}.

\begin{quote}
\textbf{\zh{我们开了一个晚会欢送他}} \emph{wǒmen kāi le yí ge wǎnhuì huānsòng tā} — We held a party to see him off.
\end{quote}

The new character, in pieces: \zh{送} — nine strokes: \zh{关} with its last stroke shrunk to a dot, inside the walking piece of \zh{这}, which wraps it from the left and runs along the foot. The next lesson puts it in your hand.

\begin{cousinweb}
\textbf{\zh{欢送}} (\emph{huānsòng}) is built from \zh{欢} 'glad' and \emph{sòng} 'to send', both inherited from Middle Chinese, \emph{xwan} and \emph{suwngH} (Guangyun, in Baxter's transcription), which glosses the first 'joy' and the second 'to send away'.
\end{cousinweb}

\begin{grammarlens}[title={What you've built}]
One more word for telling what happened.
\end{grammarlens}

\subsection*{Your turn}
\begin{itemize}
\raggedright
  \item \emph{Say it:} \emph{huānsòng}
  \item \emph{Say it:} \emph{huānsòng}, once more
  \item \emph{Say it:} \emph{wǒmen kāi le yí ge wǎnhuì huānsòng tā} — We held a party to see him off.
  \item \emph{Recall:} say the Mandarin for a string of firecrackers, then the Mandarin for to take leave of someone, then say \emph{huānsòng} again
\end{itemize}

\begin{tcolorbox}[breakable,colback=teal!4,colframe=teal!35!black,title={Before you move on}]
What does \zh{欢送} mean? (\textquotedblleft{}To give a send-off\textquotedblright{}.) Say it once more.
\end{tcolorbox}
\section[\texorpdfstring{\zh{送} (sòng)}{送 (sòng)}]{\zh{送} — \zh{关} in the walking piece}

\label{lesson:ZH-W336-song-send}



\begin{quote}
The word first, then the character inside it.

\begin{itemize}
\raggedright
  \item \emph{Recall:} say \emph{huānsòng} — \textbf{R1}, one lesson back, before it decays
\end{itemize}

\begin{quote}
\zh{关} \zh{这}
\end{quote}

\emph{Read it:} each one aloud, with nothing said first — \emph{guān}, \emph{zhè}
\end{quote}

\subsection*{Script — \zh{关} in the walking piece}
\hlblockfigure{\detokenize{figures/ZH-W336-song-send-filmstrip.pdf}}{How \zh{送} is written, stroke by stroke}

\begin{quote}
\zh{送}
\end{quote}

\textbf{Two pieces: the inside first, then the piece that wraps it.}

\textbf{The inside piece, \zh{关}, six strokes,} as you wrote it: two short strokes on top, two bars, a long stroke falling to the left, and last the stroke falling to the right, which here shrinks to a dot.

\textbf{The walking piece, three strokes,} as in \zh{这}: a dot, a small stroke that zigzags down, then a long sweep along the foot to the right.

6 + 3 = nine.

\subsection*{Your turn}
\begin{itemize}
\raggedright
  \item \emph{Write it:} \zh{送}, all nine, with the model in front of you
  \item \emph{Cover:} write it again from memory, and count the pieces, not the strokes
\end{itemize}

\begin{tcolorbox}[breakable,colback=teal!4,colframe=teal!35!black,title={Before you move on}]
How many strokes in \zh{送}? (\textbf{Nine}.) What sits inside? (\textbf{\zh{关}}.) What is written last? (\textbf{The long sweep along the foot}.)

Source: \href{https://raw.githubusercontent.com/chanind/hanzi-writer-data/68d10a4b21150cae5e1ebbd223eed289cf32d90c/data/\%E9\%80\%81.json}{Hanzi Writer Data \zh{送}, pinned stroke paths and medians}.
\end{tcolorbox}
\section[\texorpdfstring{\zh{送客} (sòngkè)}{送客 (sòngkè)}]{\zh{送客} (sòngkè) — to see a guest out}

\label{lesson:ZH-C336-songke}



\begin{quote}
Before the new one: say the Mandarin for to take leave of someone, then the Mandarin for to give a send-off.
\end{quote}

\subsection*{You'll want to know: \zh{送客}}
\textbf{\zh{送客}} — \emph{sòngkè} — \textquotedblleft{}to see a guest out\textquotedblright{}.

\begin{quote}
\textbf{\zh{客人走的时候他去门口送客}} \emph{kèrén zǒu de shíhou tā qù ménkǒu sòngkè} — When the guests left, he saw them to the door.
\end{quote}

\begin{cousinweb}
\textbf{\zh{送客}} (\emph{sòngkè}) is built from \zh{送} 'to send' and \zh{客} 'a guest', both inherited from Middle Chinese, \emph{suwngH} and \emph{khaek} (Guangyun, in Baxter's transcription), which glosses the second 'a guest'.
\end{cousinweb}

\begin{grammarlens}[title={What you've built}]
One more word for telling what happened.
\end{grammarlens}

\subsection*{Your turn}
\begin{itemize}
\raggedright
  \item \emph{Say it:} \emph{sòngkè}
  \item \emph{Say it:} \emph{sòngkè}, once more
  \item \emph{Say it:} \emph{kèrén zǒu de shíhou tā qù ménkǒu sòngkè} — When the guests left, he saw them to the door.
  \item \emph{Recall:} say the Mandarin for to take leave of someone, then the Mandarin for to give a send-off, then say \emph{sòngkè} again
\end{itemize}

\begin{tcolorbox}[breakable,colback=teal!4,colframe=teal!35!black,title={Before you move on}]
What does \zh{送客} mean? (\textquotedblleft{}To see a guest out\textquotedblright{}.) Say it once more.
\end{tcolorbox}
\section[sònglǐ]{sònglǐ — to give a present}

\label{lesson:ZH-C336-songli}



\begin{quote}
Before the new one: say the Mandarin for to give a send-off, then the Mandarin for to see a guest out.
\end{quote}

\subsection*{You'll want to know: sònglǐ}
\emph{sònglǐ} — \textquotedblleft{}to give a present\textquotedblright{}.

\begin{quote}
\emph{Chūnjié de shíhou Zhōngguó rén chángcháng sònglǐ} — At the lunar new year Chinese people often give presents.
\end{quote}

Its characters come later; for now it is a word for the ear.

\begin{cousinweb}
\emph{sònglǐ} is built from \zh{送} 'to send' and \emph{lǐ} 'a courtesy', both inherited from Middle Chinese, \emph{suwngH} and \emph{lejX} (Guangyun, in Baxter's transcription), which says of the second, quoting an older dictionary, that it is how one serves the spirits and brings blessing.
\end{cousinweb}

\begin{grammarlens}[title={What you've built}]
One more word for telling what happened.
\end{grammarlens}

\subsection*{Your turn}
\begin{itemize}
\raggedright
  \item \emph{Say it:} \emph{sònglǐ}
  \item \emph{Say it:} \emph{sònglǐ}, once more
  \item \emph{Say it:} \emph{Chūnjié de shíhou Zhōngguó rén chángcháng sònglǐ} — At the lunar new year Chinese people often give presents.
  \item \emph{Recall:} say the Mandarin for to give a send-off, then the Mandarin for to see a guest out, then say \emph{sònglǐ} again
\end{itemize}

\begin{tcolorbox}[breakable,colback=teal!4,colframe=teal!35!black,title={Before you move on}]
What does sònglǐ mean? (\textquotedblleft{}To give a present\textquotedblright{}.) Say it once more.
\end{tcolorbox}
\section[wòshǒu]{wòshǒu — to shake hands}

\label{lesson:ZH-C336-woshou}



\begin{quote}
Before the new one: say the Mandarin for to see a guest out, then the Mandarin for to give a present.
\end{quote}

\subsection*{You'll want to know: wòshǒu}
\emph{wòshǒu} — \textquotedblleft{}to shake hands\textquotedblright{}.

\begin{quote}
\emph{gàobié de shíhou wǒmen wòshǒu le} — When we said goodbye, we shook hands.
\end{quote}

Its characters come later; for now it is a word for the ear.

\begin{cousinweb}
\emph{wòshǒu} is built from \emph{wò} 'to grasp' and \zh{手} 'a hand', both inherited from Middle Chinese, \emph{'aewk} and \emph{syuwX} (Guangyun, in Baxter's transcription), which glosses the first 'to hold'.
\end{cousinweb}

\begin{grammarlens}[title={What you've built}]
One more word for telling what happened.
\end{grammarlens}

\subsection*{Your turn}
\begin{itemize}
\raggedright
  \item \emph{Say it:} \emph{wòshǒu}
  \item \emph{Say it:} \emph{wòshǒu}, once more
  \item \emph{Say it:} \emph{gàobié de shíhou wǒmen wòshǒu le} — When we said goodbye, we shook hands.
  \item \emph{Recall:} say the Mandarin for to see a guest out, then the Mandarin for to give a present, then say \emph{wòshǒu} again
\end{itemize}

\begin{tcolorbox}[breakable,colback=teal!4,colframe=teal!35!black,title={Before you move on}]
What does wòshǒu mean? (\textquotedblleft{}To shake hands\textquotedblright{}.) Say it once more.
\end{tcolorbox}
